\section{EvoDiff 的评估：质量、多样性与功能验证}

\subsection{评估方法论}

评估一个蛋白质生成模型远比评估图像生成模型复杂。Amini 强调，需要从\textbf{多个层面}审慎设计评估策略：

\begin{table}[H]
\centering
\caption{蛋白质生成模型的多层次评估框架}
\begin{tabular}{@{}lll@{}}
\toprule
\textbf{评估维度} & \textbf{核心问题} & \textbf{评估方法} \\
\midrule
个体质量 & 单个蛋白质是否生物学可行？ & 结构预测 + self-consistency \\
分布覆盖 & 生成样本是否覆盖数据分布？ & 特征空间分布可视化与比较 \\
功能实现 & 生成的蛋白质能否实现目标功能？ & 实验室湿实验验证 \\
\bottomrule
\end{tabular}
\end{table}

\begin{warningbox}{不能仅看"准确率"}
蛋白质生成模型的评估\textbf{不能简化为单一的准确率指标}。一个模型可能生成高质量的个体蛋白质，但如果所有生成的蛋白质都集中在功能空间的一个狭窄区域，那它作为设计工具的价值就大打折扣。因此，\textbf{分布层面的多样性}和\textbf{功能层面的可验证性}与个体质量同等重要。
\end{warningbox}

\subsection{个体质量评估：Self-Consistency}

EvoDiff 采用了一种巧妙的 self-consistency 评估方法：

\begin{enumerate}
    \item 用 EvoDiff 生成一条新蛋白质序列 $\mathbf{s}_{\text{gen}}$
    \item 用结构预测工具（如 AlphaFold）预测 $\mathbf{s}_{\text{gen}}$ 的三维结构 $\mathcal{S}$
    \item 从结构 $\mathcal{S}$ 出发，推断一条能折叠为该结构的序列 $\mathbf{s}_{\text{inv}}$
    \item 计算 $\mathbf{s}_{\text{gen}}$ 与 $\mathbf{s}_{\text{inv}}$ 之间的一致性
\end{enumerate}

如果一致性高，说明生成的序列能够稳定地折叠为一个确定的三维结构，是\textbf{结构上可行的} (structurally realistic and consistent)。结果显示，EvoDiff 的生成蛋白质在结构一致性方面表现良好。

\subsection{分布覆盖评估：与天然蛋白质的比较}

为评估 EvoDiff 生成样本的多样性，团队进行了分布层面的分析：

\begin{enumerate}
    \item 取天然蛋白质测试集序列，用预训练模型提取特征向量
    \item 将特征向量通过降维方法（如 t-SNE 或 UMAP）投影到 2D 空间，得到天然蛋白质的分布
    \item 对 EvoDiff 生成的大量蛋白质（数千条）执行相同操作
    \item 叠加两个分布，评估覆盖程度
\end{enumerate}

\begin{importantbox}{EvoDiff 的分布覆盖优势}
与多种基线方法的对比显示了 EvoDiff 的显著优势：
\begin{itemize}
    \item \textbf{vs. Next Token Prediction 模型}：两者覆盖范围大致相当，但 next token 模型略优
    \item \textbf{vs. Masked 一步生成模型}：EvoDiff 的覆盖范围明显更广，说明迭代去噪优于一步生成
    \item \textbf{vs. 纯结构方法}（如 RFdiffusion 等）：结构方法展现出\textbf{严重的偏差}，大量样本集中在功能空间的狭窄区域
\end{itemize}
\end{importantbox}

纯结构方法之所以存在严重偏差，根源在于结构数据的偏斜分布：实验解析的约 30 万个蛋白质结构中，球状蛋白和 $\alpha$-helix 丰富的蛋白质被过度代表。基于这些数据训练的模型自然会过采样 (oversample) 这些类型的蛋白质。而 EvoDiff 基于序列数据训练，数据量更大、覆盖更广，因此能更好地捕获天然蛋白质功能空间的多样性。

\subsection{实验验证：从计算到实验室}

最终的"黄金标准"评估是\textbf{在实验室中进行湿实验验证}。团队从 EvoDiff 的生成结果中采样候选蛋白质，进行了以下实验：

\textbf{结构稳定性验证}：
\begin{itemize}
    \item 将完全人工设计的蛋白质序列进行\textbf{生物表达}（在细胞中合成蛋白质）
    \item 测量蛋白质的结构稳定性
    \item 结果：成功表达并验证了 4 个结构稳定的全新蛋白质设计
\end{itemize}

\begin{knowledgebox}{从计算设计到实验验证的完整闭环}
AI 蛋白质设计的最终价值必须通过实验验证来证明。这涉及一个复杂的流程：(1) 计算模型生成候选序列；(2) 通过基因合成获得对应的 DNA；(3) 在细胞系统中表达蛋白质；(4) 纯化蛋白质；(5) 用物理化学手段（如圆二色谱、X 射线晶体学等）验证结构和功能。EvoDiff 团队完成了这一完整闭环，证明了模型的实际可用性。
\end{knowledgebox}

\subsection{本章小结}

EvoDiff 的评估覆盖了个体质量、分布多样性和实验验证三个层次。在分布覆盖方面，EvoDiff 相较于纯结构方法展现出明显优势，这归因于序列数据在规模和多样性上的内在优势。最重要的是，EvoDiff 生成的蛋白质在实验室中被成功表达和结构验证，证明了模型的实际可用性。


